\documentclass[11pt,a4paper]{article}
\usepackage[margin=2.4cm]{geometry}
\usepackage{graphicx}
\usepackage{caption}
\usepackage{subcaption}
\usepackage{float}
\usepackage{booktabs}
\usepackage{amsmath,amssymb}
\usepackage{xcolor}
\usepackage{listings}
\usepackage[colorlinks=true,linkcolor=blue!60!black,citecolor=blue!60!black,urlcolor=blue!60!black]{hyperref}

\lstset{
  language=Python,
  basicstyle=\ttfamily\scriptsize,
  keywordstyle=\color{blue!70!black}\bfseries,
  commentstyle=\color{green!45!black}\itshape,
  stringstyle=\color{orange!75!black},
  showstringspaces=false,
  numbers=left,
  numberstyle=\tiny\color{gray},
  frame=single,
  framerule=0.3pt,
  breaklines=true,
  tabsize=4,
}

\newcommand{\code}[1]{\texttt{\small #1}}

\title{\textbf{Assignment 1: Perceptron for Classification and Regression}\\[1em]
\large A Mechanism-Level Analysis with Full Dataset Characterisation\\[2em]
\normalsize\textbf{Group 11} --- CS601T Deep Learning\\[1.5em]
\normalsize
CS24BT055\enspace Himanshu Nagwanshi\\
CS24BT061\enspace Sidharth B\\
CS24BT060\enspace Anshuman Rahul}
\author{}
\date{\today}

\begin{document}
\maketitle
\thispagestyle{empty}

\clearpage
\setcounter{page}{1}

\begin{abstract}
This report documents the design, implementation, and empirical behaviour of a single-layer
perceptron trained with batch gradient descent, applied to two classification tasks (linearly
separable; nonlinearly separable) and two regression tasks (univariate; bivariate). We provide
a complete characterisation of all four raw datasets (feature ranges, class distributions,
correlations, geometric structure) before any processing. Beyond reporting metrics, we explain
\emph{why} each result occurs through the geometry of the decision boundary, the algebra of
the gradient update, the effect of the activation derivative on learning speed, and the
representational ceiling that causes failure on nonlinearly separable data. We cite the
academic sources for every design decision. All numbers are from actual runs of this
codebase (seed = 42).
\end{abstract}

%% ================================================================
%% SECTION 1: DATASET CHARACTERISATION
%% ================================================================
\section{Dataset Characterisation (Raw, Unprocessed)}

Before any training or preprocessing, we characterise each dataset in full. All statistics
below are computed on the \emph{complete raw files} before the 70/30 train--test split.

\subsection{Dataset 1: Linearly Separable (LS)}

\begin{figure}[H]
\centering
\includegraphics[width=0.75\linewidth]{outputs/datasets/ls_raw.png}
\caption{Raw scatter of Dataset 1. Three Gaussian clouds in $\mathbb{R}^2$, one per class.}
\label{fig:ls-raw}
\end{figure}

Three \code{.txt} files loaded by \code{load\_ls\_data()}, each containing $n \times 2$ columns
(x1, x2). Labels assigned by file sort order (Class1=0, Class2=1, Class3=2).

\begin{table}[H]
\centering
\small
\begin{tabular}{lrrrrr}
\toprule
Class & $n$ & x1 range & x1 mean $\pm$ std & x2 range & x2 mean $\pm$ std \\
\midrule
0 & 500 & $[-1.63, 6.76]$  & $3.37 \pm 1.51$ & $[-12.56, 2.37]$  & $-4.58 \pm 1.52$ \\
1 & 500 & $[-0.09, 22.12]$ & $15.18 \pm 2.75$ & $[-7.48, 18.66]$  & $11.16 \pm 2.70$ \\
2 & 500 & $[-14.62, 4.40]$ & $-5.13 \pm 2.63$ & $[-15.17, 2.79]$  & $-4.98 \pm 2.66$ \\
\midrule
\textbf{All} & \textbf{1500} & $[-14.62, 22.12]$ & $4.47 \pm 8.76$ & $[-15.17, 18.66]$ & $0.53 \pm 9.63$ \\
\bottomrule
\end{tabular}
\caption{Feature ranges and statistics for Dataset 1.}
\label{tab:ls-stats}
\end{table}

\textbf{Key observations.}
(i) Overall Pearson correlation between x1 and x2 is $-0.1173$.
(ii) Three clusters are visually separated with no overlap, confirming linear separability.
(iii) Class 1 is offset to the top-right ($x_1 \approx 15$, $x_2 \approx 11$); Classes 0 and 2
share similar $x_2$ centres ($\approx -5$) but separate along $x_1$.
(iv) Feature magnitudes reach $\sim 22$, which is relevant for understanding why the
numerically-stable logistic function is necessary (see Section~\ref{sec:logistic-stability}).

\subsection{Dataset 2: Nonlinearly Separable (NLS)}

\begin{figure}[H]
\centering
\includegraphics[width=0.75\linewidth]{outputs/datasets/nls_raw.png}
\caption{Raw scatter of Dataset 2. Three concentric crescent/ring structures.}
\label{fig:nls-raw}
\end{figure}

Loaded from a single \code{.txt} file (\code{skiprows=1}). Classes are determined by
\emph{position in file}: first 500 = class 0, next 500 = class 1, last 700 = class 2.

\begin{table}[H]
\centering
\small
\begin{tabular}{lrrrrr}
\toprule
Class & $n$ & x1 range & x1 mean $\pm$ std & x2 range & x2 mean $\pm$ std \\
\midrule
0 (inner) & 500 & $[-1.22, 0.74]$  & $-0.01 \pm 0.32$ & $[-1.24, 0.02]$  & $-0.34 \pm 0.31$ \\
1 (middle) & 500 & $[-1.86, 1.91]$  & $0.12 \pm 0.80$  & $[-1.56, 1.64]$  & $0.12 \pm 0.79$ \\
2 (outer) & 700 & $[-2.90, 2.87]$  & $0.02 \pm 1.55$  & $[-2.75, 2.92]$  & $-0.03 \pm 1.52$ \\
\midrule
\textbf{All} & \textbf{1700} & $[-2.90, 2.90]$ & $0.03 \pm 1.17$ & $[-2.75, 2.92]$ & $-0.05 \pm 1.15$ \\
\bottomrule
\end{tabular}
\caption{Feature ranges for Dataset 2. Feature magnitudes are small ($[-3,3]$) compared to LS.}
\label{tab:nls-stats}
\end{table}

\begin{table}[H]
\centering
\small
\begin{tabular}{lrrr}
\toprule
Class & $n$ & Mean dist.\ from origin & Std dist. \\
\midrule
0 (inner)  & 500 & 0.47  & 0.16 \\
1 (middle) & 500 & 1.13  & 0.22 \\
2 (outer)  & 700 & 2.17  & 0.46 \\
\bottomrule
\end{tabular}
\caption{Radial distance from origin for each class, confirming concentric ring structure.}
\label{tab:nls-radial}
\end{table}

\textbf{Key observations.}
(i) Class sizes are \emph{imbalanced}: 500, 500, 700 --- stratified splitting is essential.
(ii) The concentric ring structure means no linear boundary can separate any two class pair: every pair requires a non-linear surface (e.g.\ $x_1^2 + x_2^2 = c$).

\subsection{Dataset 3: Univariate Regression}

\begin{figure}[H]
\centering
\includegraphics[width=0.75\linewidth]{outputs/datasets/univariate_raw.png}
\caption{Raw scatter of Dataset 3. Non-linear periodic function.}
\label{fig:uni-raw}
\end{figure}

\begin{table}[H]
\centering
\begin{tabular}{lrr}
\toprule
Variable & Range & Mean $\pm$ std \\
\midrule
$x$ & $[0.00, 9.99]$   & $5.03 \pm 2.89$ \\
$y$ & $[-3.96, 5.96]$   & $0.90 \pm 2.56$ \\
\midrule
$n$ & 1000 & \\
\bottomrule
\end{tabular}
\caption{Univariate dataset statistics.}
\label{tab:uni-stats}
\end{table}

The target function is a noisy periodic signal on $x \in [0,10]$. The mean target
$\bar{y} = 0.90$ is small relative to the signal amplitude ($\sim 10$ units peak-to-peak),
which inflates \%RMSE values.

\subsection{Dataset 4: Bivariate Regression}

\begin{figure}[H]
\centering
\includegraphics[width=0.65\linewidth]{outputs/datasets/bivariate_raw.png}
\caption{Raw 3-D scatter of Dataset 4. The target surface is visibly curved.}
\label{fig:biv-raw}
\end{figure}

\begin{table}[H]
\centering
\begin{tabular}{lrr}
\toprule
Variable & Range & Mean $\pm$ std \\
\midrule
$x_1$ & $[0.00, 10.00]$ & $5.01 \pm 2.88$ \\
$x_2$ & $[0.00, 10.00]$ & $4.98 \pm 2.90$ \\
$y$   & $[-1.00, 27.02]$ & $12.66 \pm 6.49$ \\
\midrule
$n$ & 10201 & \\
\bottomrule
\end{tabular}
\caption{Bivariate dataset statistics. $n = 10{,}201$ samples on an approximate $101 \times 101$ grid.}
\label{tab:biv-stats}
\end{table}

%% ================================================================
%% SECTION 2: MODEL ARCHITECTURE
%% ================================================================
\section{Model Architecture}

\subsection{The Single-Layer Perceptron}

The core unit is a single neuron with an affine projection followed by a scalar non-linearity
(\code{models/perceptron.py}):
\begin{equation}
\hat{y} = \phi(\mathbf{x}^{\top}\mathbf{w} + b), \qquad
\phi \in \{\sigma, \tanh, \mathrm{id}\}.
\label{eq:forward}
\end{equation}

\begin{itemize}
\item \textbf{Logistic} $\sigma(z) = 1/(1+e^{-z})$: outputs $(0,1)$, thresholded at $0.5$.
\item \textbf{Tanh}: outputs $(-1,1)$, thresholded at $0$. Zero-centred, which speeds
convergence because its mean activation is closer to zero \cite{lecun1998efficient}.
\item \textbf{Linear} $\phi(z)=z$: used for regression; $\phi'=1$ always.
\end{itemize}

\subsection{Numerically-Stable Logistic Function}\label{sec:logistic-stability}

Implementation in \code{perceptron.py} (line 10):
\begin{lstlisting}
def sigmoid(self, x):
    return np.where(x >= 0, 1/(1+np.exp(-x)),
                          np.exp(x)/(1+np.exp(x)))
\end{lstlisting}

For large negative $z$, $\exp(-z)$ overflows \code{float64}. The split form keeps every
exponent argument $\le 0$. This is the standard engineering practice in TensorFlow
\code{tf.nn.sigmoid} \cite{numpy2024}. Necessary because Dataset 1 features reach $\sim 22$,
pushing pre-activations past the overflow threshold.

\subsection{Gradient Update Rule}\label{sec:gradient}

MSE loss $\mathcal{L} = \frac{1}{N}\sum_i (\hat y_i - y_i)^2$. Chain rule gives
(\code{perceptron.py}, lines 26--30):
\begin{equation}
\delta_i = (\hat y_i - y_i)\,\phi'(\hat y_i), \qquad
\nabla_{\mathbf{w}} \mathcal{L} = \frac{1}{N}\mathbf{X}^{\top}\boldsymbol{\delta}, \qquad
\nabla_b \mathcal{L} = \frac{1}{N}\sum_i \delta_i .
\end{equation}

Peak derivative magnitudes: logistic $\max = 0.25$, tanh $\max = 1.0$ ($4\times$ larger),
linear $\phi' = 1$ always.

\subsection{Optimiser and Initialisation}

Batch GD, $\eta = 0.01$, 100 epochs. Chosen over SGD because with $N \le 10{,}201$ and
$d \le 2$, full-batch steps cost microseconds and produce smooth deterministic curves.
$\eta = 0.01$ empirically chosen: larger ($\ge 0.5$) saturates the logistic function; smaller ($\le 10^{-3}$)
leaves NLS curves flat. Weights $\mathcal{N}(0, 0.01^2)$, $b=0$ \cite{lecun1998efficient}.

%% ================================================================
%% SECTION 3: MULTI-CLASS STRATEGY
%% ================================================================
\section{Multi-Class Strategy: One-vs-One (OvO)}

For 3-class problems, \code{models/one\_vs\_one.py} trains $\binom{3}{2}=3$ independent
binary perceptrons, one per class pair \cite{hastie1998pairwise}.

\subsection{Why OvO and Not One-vs-Rest?}
\begin{itemize}
\item \textbf{Balanced sub-problems.} Each OvO classifier sees 1000 samples split evenly.
OvR would give 500:1000 imbalance.
\item \textbf{Simpler task.} Each binary problem is genuinely separable on Dataset 1.
\item \textbf{Cost is trivial.} With $k=3$, OvO's balance advantage dominates.
\end{itemize}

\subsection{Label Encoding per Activation}

Logistic uses $\{0,1\}$ labels (threshold at $0.5$); tanh uses $\{-1,+1\}$ (threshold at $0$).
This is critical: using $\{0,1\}$ with tanh pushes all predictions positive and breaks voting.

\subsection{Voting}
Each pairwise classifier casts one vote; \code{argmax} determines the final class.
Ties broken by lowest class index.

%% ================================================================
%% SECTION 4: CLASSIFICATION RESULTS
%% ================================================================
\section{Classification Results}\label{sec:ls-results}

Seed $= 42$, $\eta = 0.01$, 100 epochs, 70/30 stratified split.

\subsection{Dataset 1: Linearly Separable}

\begin{table}[H]
\centering
\begin{tabular}{lcccc}
\toprule
Activation & Accuracy & Precision & Recall & F1 \\
\midrule
Logistic & 1.000 & 1.000 & 1.000 & 1.000 \\
Tanh     & 1.000 & 1.000 & 1.000 & 1.000 \\
\bottomrule
\end{tabular}
\caption{LS test-set metrics ($n_{\text{test}}=450$, 150/class).}
\label{tab:ls-metrics}
\end{table}

\begin{table}[H]
\centering
\begin{tabular}{lccc}
\toprule
Pair & Logistic MSE & Tanh MSE \\
\midrule
(0 vs 1) & 0.0118 & 0.0032 \\
(0 vs 2) & 0.0270 & 0.0081 \\
(1 vs 2) & 0.0217 & 0.0083 \\
\bottomrule
\end{tabular}
\caption{Final MSE per OvO classifier. Tanh converges $2$--$3\times$ faster.}
\label{tab:ls-mse}
\end{table}

\textbf{Why it works.} Each class pair admits a separating line. A perceptron implements
a half-plane; MSE is convex through the monotone activation. Tanh's lower MSE comes from
$\tanh'(0)=1$ vs $\sigma'(0)=0.25$.

\begin{figure}[H]
\centering
\includegraphics[width=0.85\linewidth]{outputs/ls/error_curves_logistic.png}
\caption{Error vs.\ epochs, LS, logistic activation. Smooth monotone convergence.}
\label{fig:ls-error}
\end{figure}

\begin{figure}[H]
\centering
\includegraphics[width=0.48\linewidth]{outputs/ls/decision_regions_logistic.png}
\hfill
\includegraphics[width=0.48\linewidth]{outputs/ls/decision_regions_tanh.png}
\caption{Decision regions: logistic (left) vs tanh (right). Both produce clean partitions.}
\label{fig:ls-dr}
\end{figure}

\subsection{Dataset 2: Nonlinearly Separable}\label{sec:nls-failure}

\begin{table}[H]
\centering
\begin{tabular}{lcccc}
\toprule
Activation & Accuracy & Prec. & Rec. & F1 \\
\midrule
Logistic & 0.266 & 0.115 & 0.215 & 0.149 \\
Tanh     & 0.280 & 0.112 & 0.226 & 0.150 \\
Chance   & 0.333 & ---     & ---     & ---  \\
\bottomrule
\end{tabular}
\caption{NLS test-set ($n_{\text{test}}=511$). Both at or below majority-class chance.}
\label{tab:nls-metrics}
\end{table}

\begin{table}[H]
\centering
\begin{tabular}{lccc}
\toprule
Pair & Logistic MSE & Tanh MSE \\
\midrule
(0 vs 1) & 0.2489 & 0.9554 \\
(0 vs 2) & 0.2483 & 0.9555 \\
(1 vs 2) & 0.2491 & 0.9755 \\
\bottomrule
\end{tabular}
\caption{NLS final MSE. Logistic stalls at chance ($\approx 0.25$); tanh diverges ($\approx 0.96$).}
\label{tab:nls-mse}
\end{table}

\begin{table}[H]
\centering
\begin{tabular}{lrrr}
\toprule
\multicolumn{4}{c}{\textbf{Logistic Confusion Matrix}} \\
\midrule
 & Pred 0 & Pred 1 & Pred 2 \\
Actual 0 & 0 & 0 & 150 \\
Actual 1 & 40 & 0 & 110 \\
Actual 2 & 75 & 0 & 136 \\
\bottomrule
\end{tabular}
\caption{All predictions collapse to class 2 (the training majority).}
\label{tab:nls-cm}
\end{table}

\textbf{Why it fails---representation, not optimisation.} Concentric rings require
a non-linear boundary ($x_1^2 + x_2^2 = c$). A half-plane cannot isolate an annulus.
Logistic stalls at MSE $\approx 0.25$ (balanced binary chance level). Tanh diverges
(MSE $> 0.95$) because its unbounded range allows predictions to drift to $\pm 1$
with equal confidence, collapsing the voting scheme.

\begin{figure}[H]
\centering
\includegraphics[width=0.85\linewidth]{outputs/nls/error_curves_logistic.png}
\caption{Error vs.\ epochs, NLS, logistic activation. Flat at $\approx 0.25$---chance-level plateau.}
\label{fig:nls-error}
\end{figure}

\begin{figure}[H]
\centering
\includegraphics[width=0.48\linewidth]{outputs/nls/decision_regions_logistic.png}
\hfill
\includegraphics[width=0.48\linewidth]{outputs/nls/decision_regions_tanh.png}
\caption{NLS decision regions: logistic (left) vs tanh (right). Concentric structure unresolvable.}
\label{fig:nls-dr}
\end{figure}

%% ================================================================
%% SECTION 5: REGRESSION RESULTS
%% ================================================================
\section{Regression Results}

A single perceptron with identity activation implements $f(\mathbf{x})=\mathbf{x}^{\top}\mathbf{w}+b$;
trained by GD on MSE it converges to the best affine approximation.

\begin{table}[H]
\centering
\begin{tabular}{lccccc}
\toprule
Task & $n$ (train/test) & Train RMSE & \%RMSE & Test RMSE & \%RMSE \\
\midrule
Univariate ($d=1$) & 700/301 & 1.2842 & 50.22 & 1.2662 & 53.52 \\
Bivariate ($d=2$) & 7140/3061 & 3.3674 & 82.53 & 3.3342 & 83.64 \\
\bottomrule
\end{tabular}
\caption{Regression performance. \%RMSE $=$ RMSE$/\bar{y} \times 100$.}
\label{tab:reg}
\end{table}

\textbf{Underfitting diagnosis.} Train and test RMSE are nearly identical (1.28 vs 1.27;
3.37 vs 3.33)---evidence of \emph{high bias} (underfitting), not overfitting. An affine
model captures the linear component of the target but cannot represent the sinusoidal
structure (univariate) or the curved surface (bivariate).

The large \%RMSE values for the univariate task are partly an artefact of the normalisation
divisor $\bar{y} = 0.90$ being small relative to the signal amplitude ($\sim 10$ units
peak-to-peak). The absolute RMSE of 1.27 on a $y$-range of $[-4, 6]$ corresponds to
$\sim 13\%$ of the total signal range.

\begin{figure}[H]
\centering
\includegraphics[width=0.48\linewidth]{outputs/regression/univariate/model_vs_target_test.png}
\hfill
\includegraphics[width=0.48\linewidth]{outputs/regression/univariate/scatter_test.png}
\caption{Univariate: model fit on test data (left) and target vs predicted scatter (right).}
\label{fig:reg-uni}
\end{figure}

%% ================================================================
%% SECTION 6: DESIGN DECISIONS AND ALTERNATIVES
%% ================================================================
\section{Design Decisions and Alternatives Considered}

\subsection{Feature Scaling}
Dataset 1 spans $[-15, 22]$. Unnormalised inputs push pre-activations into the logistic
function's saturated tails. We chose not to rescale, keeping the pipeline faithful to raw
data. Standardising to zero-mean/unit-variance would speed convergence; omitted for simplicity.

\subsection{Loss: MSE vs Cross-Entropy}
The assignment specifies MSE. Cross-entropy with logistic output gives $\delta = \hat y - y$ (no
$\sigma'$ factor), producing larger gradients and faster convergence. We follow the spec.

\subsection{Tanh Label Encoding}
Using $\{0,1\}$ with tanh would place the optimum at $\hat y = 0.5$, misaligning with
tanh's $[-1,1]$ range. The $\{-1,+1\}$ encoding is the standard fix.

\subsection{Stratified Split}
Essential because NLS has imbalanced classes (500, 500, 700). Without stratification,
class 2 could be underrepresented in training.

%% ================================================================
%% SECTION 7: REFERENCES
%% ================================================================
\begin{thebibliography}{99}

\bibitem{lecun1998efficient}
Y.~LeCun, L.~Bottou, G.~B.~Orr, and K.-R.~M{\"u}ller,
``Efficient BackProp,''
\emph{Neural Networks: Tricks of the Trade}, Springer, 1998.
Available: \url{https://leon.bottou.org/publications/pdf/tricks-2012.pdf}

\bibitem{hastie1998pairwise}
T.~Hastie and R.~Tibshirani,
``Classification by Pairwise Coupling,''
\emph{Annals of Statistics}, vol.~26, no.~2, pp.~451--471, 1998.
DOI: \url{https://doi.org/10.1214/aos/1028144843}

\bibitem{numpy2024}
NumPy Developers,
``NumPy v2.1 Documentation,''
\url{https://numpy.org/doc/stable/}

\bibitem{goodfellow2016deep}
I.~Goodfellow, Y.~Bengio, and A.~Courville,
\emph{Deep Learning}, MIT Press, 2016.
\url{https://www.deeplearningbook.org/}

\bibitem{haykin2009neural}
S.~Haykin,
\emph{Neural Networks and Learning Machines}, 3rd ed., Pearson, 2009.

\end{thebibliography}

\end{document}
